% V2第24章思考题24.1--24.11独立解，原图24.27/24.28/24.37需人工重绘复核。
\section*{第24章\quad 光的偏振：思考题解析}
\paragraph{思考题 24.1}
线偏振、椭圆偏振光是两相互正交电场分量
有\emph{稳定相对相位差}的相干合成，
例如
\[
\boldsymbol E(t)=
E_x\cos\omega t\,\hat{\boldsymbol x}
+E_y\cos(\omega t+\delta)\,\hat{\boldsymbol y}.
\]
若 \(\delta\) 恒定，电矢量末端描绘固定直线
或椭圆。自然光中大量独立发光微源的方向
和相位随机变化，时间或统计平均的相干矩阵
非对角交叉项为零：
\(\langle E_xE_y^*\rangle=0\)。
因此可以用两个正交分量\emph{分解能流}，
但不能将其当作一对固定相位的单色平面波，
这正是自然光与确定椭圆偏振光的区别。

\paragraph{思考题 24.2}
旋转一个理想线偏振分析器：
完全线偏振入射时
\[
\boxed{I(\theta)=I_0\cos^2(\theta-\theta_0)},
\]
会出现真正\emph{消光}角。
部分线偏振的强度随角度调制但
\(I_{\min}>0\)；
自然无偏振光强度恒为
\(\boxed{I_0/2}\)。
应注意圆偏振光也在转动线偏振器时
给出常强度，但\emph{不等于自然光}；
若圆偏振也列为待判选项，则须加四分之一波片。

\paragraph{思考题 24.3}
用已知振动方向的线偏振光穿过未知偏振片，
缓慢旋转并记录透过光强。
透过光达到最大时，未知片透振轴
与已知电矢量方向平行；
最小/消光时两者互相垂直：
\[
\boxed{I=I_0\cos^2\alpha}.
\]
仅看未标注的偏振片外观不能知道透振轴；
若只有自然光和这一片，则旋转它强度不变，
必须另用有方向参考的线偏振源或第二片。

\paragraph{思考题 24.4}
对于两个均匀透明、无吸收的介质界面，
若折射率不同，且入射光有p偏振分量，
在入射角恰为
\[
\boxed{\theta_B=\arctan(n_2/n_1)}
\]
时，p偏振的Fresnel反射振幅
\(r_p=0\)。若入射是纯p偏振，
理想界面可以无反射光。
但若两侧\emph{折射率相同}
\(n_1=n_2\)，理论上也可无反射，
并不唯一推出Brewster角。
实际未检出弱反射还可能由探测灵敏度、
多层膜消反射和涂层等造成，
不能仅由“无反光”断定偏振方向唯一。

\paragraph{思考题 24.5}
令 \(P_1\) 沿0°，\(P_2\) 沿90°，
中间 \(P_3\) 透振轴与前片夹 \(\theta\)。
对已通过 \(P_1\) 的线偏振光：
\[
\boxed{I_{\rm out}=I_1\cos^2\theta\sin^2\theta
=\frac{I_1}{4}\sin^2(2\theta)}.
\]
\(0<\theta<90^\circ\) 可透光，
\(\theta=45^\circ\) 时最大
\(I_{\rm out}=I_1/4\)；
\(\theta=0^\circ\) 或90°仍完全消光。
中间偏振片改变了可透分量的投影方向，
不违反能量守恒。

\paragraph{思考题 24.6}
按经典Thomson散射，X射线电场
使自由或近自由电子沿某方向
加速振荡。振荡电偶极辐射角分布
\[
\boxed{\frac{dP}{d\Omega}\propto\sin^2\vartheta},
\]
\(\vartheta\) 为观察方向与电子振荡轴的夹角。
沿振荡轴方向没有经典偶极辐射，
垂直于振荡轴方向辐射最大。
第一块散射使进入第二块的光带有
确定的横向偏振性质；
若第二块在所观察的垂直方向强度为零，
与入射电矢量方向的选择性以及电偶极子零辐射轴一致。
这支持X射线是具有\emph{横向电场分量}的电磁波。
真实实验还需讨论入射偏振、散射角和仪器几何，
并不能仅凭某一处零计数构成无需背景核验的绝对证明；
原图24.37需逐矢量重绘。

\paragraph{思考题 24.7}
双折射是否出现主要取决于\emph{晶体内}
传播波法线与光轴的夹角，而不是仅看
晶体外入射光束与晶体光轴是否平行。
若界面倾斜，
光在界面满足相位匹配/切向波矢守恒，
进入晶体后两个本征波的方向和
有效折射率一般不同，
不一定仍沿光轴，
因此\emph{可能出现双折射}。
只有特殊传播方向（如晶体内严格沿光轴）
两本征波折射率退化，
才无普通的双折射角分离：
\[
\boxed{\text{外部光线平行光轴，并不保证晶体内无双折射。}}
\]

\paragraph{思考题 24.8}
先用转动线偏振分析器：
能完全消光说明线偏振；
有非零最小值且强度调制说明非圆椭圆偏振
（或部分偏振）；强度不随角度变则可能
是圆偏振光或自然光。
再插入四分之一波片并调整其光轴：
圆偏振经四分之一波片变成线偏振，
能够被后分析器完全消光；
自然光经过理想无损四分之一波片仍无偏振，
分析器强度仍不随角度变化。
一般椭圆偏振在其主轴对齐四分之一波片时
可转成线偏振；部分偏振光通常仍含不能消光的
随机分量。
\[
\boxed{\text{判别圆偏振/自然光必须增加相位延迟元件。}}
\]

\paragraph{思考题 24.9}
在未知三片中逐一组合任意两片，
通过自然光并相对旋转。
两片真正的线偏振片可以在某相对角度
产生接近完全消光；
四分之一波片与线偏振片组合在
非选择偏振入射下通常不会仅凭旋转产生这样的
理想交叉双片消光。
据此识别两片偏振片后，
余下一片为四分之一波片；
再用已知线偏振片准备线偏振光，
令它与波片本征轴约45°，
输出接近圆偏振，
用后分析器旋转观察恒强度进一步验证。
\emph{非理想吸收、偏振片漏光和实际延迟量}
会使实验现象偏离理想。

\paragraph{思考题 24.10}
晶体分出互相正交的o/e本征偏振模式，
相位差
\(\delta=2\pi(n_e-n_o)d/\lambda\)。
前偏振片 \(P_1\) 把自然光转成单一稳定
线偏振矢量，确保其o/e投影彼此\emph{相干}；
后分析器 \(P_2\) 再把两正交振动
投影到同一方向，才能出现交叉强度项
\(2\Re(E_oE_e^*)\)。
\[
\boxed{\text{去掉}P_1:
\text{自然光输入通常不形成稳定相位差的可见干涉};
\quad \text{去掉}P_2:\text{o/e正交项不相干涉}.}
\]
不应将两束仅有同频和同一路径就当成
必然会直接干涉。

\paragraph{思考题 24.11}
原装置为
\(P_1\to C\to P_2\)，
先得到线偏振，再由晶体产生相干正交分量、
分析器混合。
若改为
\(\boxed{C\to P_1\to P_2}\)，
自然入射光先经过无损晶体仍为统计无偏振光，
随后 \(P_1\) 才选取一个方向，
不再保留先前晶体造成的两模相位差对最终强度的
可观察干涉。即原偏振干涉式不能直接使用。
若 \(P_1\perp P_2\) 且两者在晶体后，
理想透射强度为零；
若改成 \(P_1\parallel P_2\)，
自然光经第一偏振片后强度 \(I_0/2\)，
第二片不额外损耗：
\[
\boxed{I_{\rm out}=I_0/2}
\]
（忽略吸收），与晶体延迟相位无关。
本结论以原图为两片线偏振片和无损双折射片为前提；
原图24.28相位定义还须精确重绘核验。
